\chapter{Introduction}
\label{ch:introduction}

Contemporary software development operates in a densely instrumented 
ecosystem. Every state transition recorded in an issue tracker such as 
Jira, every commit pushed to a Git repository, every build triggered in 
a continuous integration platform, every pull request reviewed on GitHub 
generates time-stamped event data that documents who performed which 
activity, when, and in what sequence. Collectively, these events 
constitute a high-fidelity chronicle of the software development process 
as it actually occurs, rather than as it is documented in process manuals 
or assumed in retrospective meetings.

The discipline of process mining, formalized by van der 
Aalst~\cite{aalst2016process}, has demonstrated that event data can be 
systematically transformed into process models, compliance diagnostics, 
performance metrics, and predictions of future behavior. The field has 
matured along an analytical spectrum that progresses from descriptive 
capabilities (what happened?), through diagnostic (why did it happen?) 
and predictive (what will happen?), to prescriptive (what should be 
done?). In domains such as healthcare, manufacturing, financial services, 
and public administration, process mining has delivered measurable 
organizational value~\cite{aalst2012process}.

Software engineering, however, presents a paradox. Despite generating 
some of the richest and most fine-grained event data among all 
organizational activities, the application of process mining to the 
Software Development Lifecycle (SDLC) remains limited to descriptive 
analysis. Development teams continue to rely on aggregated metrics 
(velocity, throughput, cycle time) computed from project management 
dashboards, rarely interrogating the underlying process flows that 
produce these metrics. When forecasting is attempted, it typically uses 
throughput-based heuristics that treat the development process as a 
black box, ignoring the internal state dynamics, rework loops, and 
path-dependent variability that characterize real software processes.

This thesis proposes a method that bridges the gap between the rich 
event data generated by software development tools and the probabilistic 
forecasting capabilities that this data could support. The method, named 
PATHCAST integrates process mining, Markov chain estimation, Monte 
Carlo simulation, and machine learning into a unified pipeline that 
transforms heterogeneous repository data into calibrated probabilistic 
forecasts of software process outcomes.


% =============================================================================
\section{Research Problem}
\label{sec:problem}

Three interrelated problems motivate this research.

The first problem is the \textbf{descriptive ceiling} in software process 
mining. The existing literature on process mining applied to software 
development~\cite{cook1998, poncin2011process, rubin2007} 
has demonstrated that process models can be discovered from repository data. 
However, these studies stop at the descriptive level: they reveal what 
happened, but do not advance to predicting what will happen or quantifying 
the uncertainty associated with future outcomes. No published work proceeds 
from software process discovery to stochastic modeling and probabilistic 
forecasting using the discovered process structure as the foundation for 
prediction.

The second problem is the \textbf{structural blindness} of existing 
forecasting approaches. The most widely adopted forecasting method in agile 
software development, throughput-based Monte Carlo 
simulation~\cite{vacanti2015}, samples daily throughput values 
from historical data and simulates backlog completion. While practical, this 
approach treats the development process as a black box: it does not model 
internal state transitions, rework probabilities, state-specific delays, or 
path-dependent variability. An item that has been reworked three times is 
treated identically to one progressing linearly, because the simulation has 
no awareness of process structure. This limitation produces forecasts that 
cannot capture the variance explained by process dynamics.

The third problem is \textbf{technical fragmentation}. The techniques 
required for process-aware forecasting, including process mining for model 
discovery, Markov chains for transition modeling, Monte Carlo simulation 
for distributional prediction, and machine learning for case-level 
classification, exist independently in the literature. Their integration 
into a coherent pipeline, where the output of each technique serves as 
input to the next, has not been formalized or empirically validated in the 
software engineering domain. Researchers and practitioners who wish to 
combine these techniques must construct \emph{ad hoc} integrations without 
guidance from established methods.


% =============================================================================
\section{Objectives}

The growing volume of data generated throughout the Software Development Life Cycle (SDLC) offers relevant opportunities for analyzing and forecasting the behavior of software development processes. Systems such as version control repositories, \textit{issue} management platforms, and continuous integration \textit{pipelines} record in detail the activities performed during development. However, a large portion of the approaches used in practice to estimate project deadlines and outcomes still relies on aggregated metrics and heuristic methods that treat development processes as opaque systems. As a result, these approaches have limitations in capturing the structural dynamics of processes and the variability associated with different execution paths.

In this context, this research investigates the potential of process mining, combined with stochastic modeling and machine learning techniques, to improve the forecasting of software development processes.
\label{sec:objectives}

\subsection{General Objective}

To develop a computational method for the probabilistic forecasting of software development processes, based on the integration of process mining, stochastic modeling through Markov chains, Monte Carlo simulation, and machine learning, capable of transforming data from software repositories into analytical models that enable the prediction of the behavior and outcomes of these processes.

\subsection{Specific Objectives}

To achieve the general objective and test the proposed hypotheses, the following specific objectives are defined:

\begin{enumerate}
    \item Build structured event logs from data originating from software repositories, including issue tracking systems, version control repositories, and continuous integration pipelines, in order to formally represent the actual behavior of development processes.
    
    \item Apply process mining techniques to discover models that represent the actual flow of activities in software development, allowing the identification of states, transitions, and execution patterns of the process..
    
    \item Develop a probabilistic prediction model based on Markov chains and Monte Carlo simulation, capable of estimating execution time distributions, process trajectories, and possible outcomes of development processes.
    
    \item Empirically evaluate the proposed method by integrating it with machine learning techniques and comparing it with predictive models based solely on traditional software repository metrics, analyzing the gains in predictive accuracy and interpretability..
\end{enumerate}


% =============================================================================
\section{Research Questions}
\label{sec:rqs}

The research is guided by the following central question:
Can the explicit incorporation of process structure, obtained through process mining and stochastic modeling techniques, improve the accuracy and interpretability of forecasting software development processes compared to approaches based exclusively on traditional software repository metrics?


\textbf{H1 (Predictive Superiority of Process-Derived Features):} 
\emph{Predictive models that incorporate features derived from process structure achieve superior predictive performance compared with models based solely on traditional software repository metrics.}

This hypothesis evaluates whether process-aware representations provide additional explanatory and predictive power beyond conventional repository-based indicators.

\textbf{H2 (Fidelity and Interpretability of Markov-Based Stochastic Models):} 
\emph{Stochastic models based on Markov chains enable a more faithful representation of the dynamics of software development processes, allowing the generation of interpretable probabilistic forecasts.}

This hypothesis examines whether state-transition models can adequately capture process dynamics and support probabilistic predictions that are both analytically grounded and interpretable.

\textbf{H3 (Performance Gain from Multi-Technique Integration):} 
\emph{The integration of process mining, Monte Carlo simulation, and machine learning improves forecast quality when compared with the isolated use of each of these techniques.}

This hypothesis assesses whether combining complementary analytical techniques produces better forecasting performance than applying each method independently.


% =============================================================================
\section{Justification and Research Gap}
\label{sec:justification}

Several systematic reviews and surveys have been published in areas 
adjacent to the scope of this research, each covering a distinct facet 
of the landscape without addressing their convergence in the software 
engineering domain.

In the process mining literature, Augusto et al.~\cite{augusto2019} conducted a comprehensive review of 
automated process discovery algorithms. Burattin~\cite{burattin2022} 
formalized the emerging subfield of streaming process mining. Neither work 
addresses the specific challenges of applying process mining to software 
development event logs, where case correlation across heterogeneous tools 
constitutes a prerequisite not present in traditional BPM scenarios.

In the Mining Software Repositories (MSR) community, 
Hassan~\cite{hassan2008road} established a foundational research roadmap, 
and Kagdi, Collard, and Maletic~\cite{kagdi2007} produced an 
extensive taxonomy of MSR approaches. These works focus predominantly on 
mining individual artifacts (source code, commits, bug reports) rather 
than mining the process flows that connect these artifacts throughout the 
development lifecycle.

In the predictive and prescriptive monitoring literature, 
Márquez-Chamorro et al.~\cite{marquezChamorro2018} mapped 
predictive monitoring techniques and Kubrak et 
al.~\cite{kubrak2022} mapped the state of prescriptive 
process monitoring. Both surveys treat business processes generically, 
without examining the specific context of software development.

In the stochastic forecasting literature for software projects, 
Vacanti~\cite{vacanti2015} popularized throughput-based Monte 
Carlo simulation for agile delivery forecasting, while Whittaker and 
Thomason~\cite{whittaker1994} pioneered Markov chain models for 
software testing. Neither work connects these stochastic methods to 
process mining-discovered models.

Table~\ref{tab:related-surveys} synthesizes the positioning of this 
research relative to existing surveys and methods, demonstrating that 
the addressed intersection, combining process mining, stochastic 
modeling, and SDLC-specific application, remains unoccupied.

\begin{table}[htbp]
\centering
\caption{Positioning relative to existing work.}
\label{tab:related-surveys}
\footnotesize
\begin{tabular}{lccccc}
\toprule
\textbf{Work} & \textbf{PM} & \textbf{SDLC} & \textbf{Markov/} & \textbf{ML} & \textbf{Integrated} \\
 & & \textbf{Focus} & \textbf{MC Sim} & & \textbf{Pipeline} \\
\midrule
Cook \& Wolf (1998) & Yes & Discovery & No & No & No \\
Kagdi et al. (2007) & No & Artifacts & No & No & No \\
Hassan (2008) & No & Artifacts & No & No & No \\
Poncin et al. (2011) & Yes & Discovery & No & No & No \\
Whittaker \& Thomason (1994) & No & Testing & Markov & No & No \\
Vacanti (2015) & No & Forecasting & MC & No & No \\
Augusto et al. (2019) & Yes & No & No & No & No \\
Márquez-Chamorro et al. (2018) & Yes & No & No & Yes & No \\
Kubrak et al. (2022) & Yes & No & No & Yes & No \\
Rozinat et al. (2009) & Yes & No & MC & No & Partial \\
\textbf{This thesis (2026)} & \textbf{Yes} & \textbf{SDLC} & \textbf{Both} & \textbf{Yes} & \textbf{Yes} \\
\bottomrule
\end{tabular}
\end{table}

The table reveals two patterns. First, works that apply process mining to 
SDLC data stop at descriptive analysis (discovery), without advancing to 
stochastic modeling or prediction. Second, works that employ stochastic 
methods for software forecasting (Vacanti, Whittaker) operate without 
process mining, treating the process as either a black box (throughput 
sampling) or a theoretical construct (usage-model Markov chains). The 
proposed method occupies the intersection: process mining provides the 
empirical process model that parameterizes the stochastic simulation, 
applied to the SDLC domain.


% =============================================================================
\section{Expected Contributions}
\label{sec:contributions}

This thesis aims to make the following contributions:

\textbf{C1 (Method):} A formally specified method (PATHCAST) that 
integrates process mining, Markov chain estimation, Monte Carlo 
simulation, and machine learning into a four-stage pipeline for 
probabilistic forecasting of software development processes. The method 
is defined through inter-stage data contracts, formal properties, and 
correctness guarantees.

\textbf{C2 (Event Log Construction):} A configurable framework for
constructing structured event logs from heterogeneous software 
repositories, addressing the case correlation problem through a cascade 
of explicit, structural, and heuristic strategies with formal quality 
metrics. This framework is reusable independently of the forecasting 
pipeline.

\textbf{C3 (Process-Aware Stochastic Forecasting):} An empirically 
validated demonstration that Monte Carlo simulation parameterized by 
Markov chain models estimated from process mining-discovered models 
produces more accurate and better-calibrated forecasts than 
throughput-based Monte Carlo, for software development processes.

\textbf{C4 (Process-Derived Feature Set):} A novel set of features 
derived from process mining and Markov analysis (conformance scores, 
rework counts, Markov-expected remaining times) for machine learning-based 
outcome prediction in software processes, with empirical comparison 
against traditional MSR features.

\textbf{C5 (Empirical Evidence):} A reproducible empirical evaluation 
on 48 open-source software repositories, providing baseline results 
for stochastic process forecasting in the SDLC domain and a reusable 
dataset for future research.


% =============================================================================
\section{Research Methodology}
\label{sec:methodology-overview}

This research adopts the Design Science Research (DSR) paradigm, following 
the framework of Hevner et al.~\cite{hevner2004} and the methodology 
of Peffers et al.~\cite{peffers2007}. DSR is the appropriate paradigm 
because the primary goal is to design, implement, and evaluate an artifact 
(the PATHCAST method and its supporting tool) that addresses a 
practical problem identified in the literature.

The DSRM process is instantiated as follows:

\textbf{Activity 1 (Problem Identification):} The research problem is 
identified through analysis of the literature at the intersection of 
process mining, MSR, and stochastic forecasting, revealing the descriptive 
ceiling, structural blindness, and technical fragmentation described in 
Section~\ref{sec:problem}.

\textbf{Activity 2 (Objectives):} The objectives are derived from the 
identified gaps: a method that advances beyond descriptive analysis, 
incorporates process structure into forecasting, and integrates 
complementary techniques into a formally specified pipeline.

\textbf{Activity 3 (Design and Development):} The PATHCAST method 
is designed as a four-stage pipeline with formal inter-stage contracts, 
implemented as a Python-based tool using PM4Py for process mining, NumPy 
for Markov computation, and scikit-learn/XGBoost for machine learning.

\textbf{Activity 4 (Demonstration):} The method is demonstrated on 48 
open-source software repositories, producing event logs, process models, 
Markov estimates, Monte Carlo forecasts, and ML predictions.

\textbf{Activity 5 (Evaluation):} The evaluation is structured into two 
complementary stages:

\begin{itemize}
    \item \textbf{E1 (Pipeline Validation):} Assesses the quality of event 
    log construction (coverage, correlation rate) and process discovery 
    (fitness, precision, generalization) across the 48 repositories. 
    Addresses RQ1.
    
    \item \textbf{E2 (Forecasting Accuracy):} Compares the proposed 
    process-aware Monte Carlo against baselines (throughput-based MC, 
    analytical Markov, historical mean) using temporal hold-out validation. 
    Evaluates ML feature sets through cross-validated comparison. 
    Addresses RQ2 and RQ3.
\end{itemize}

\textbf{Activity 6 (Communication):} Results are communicated through 
this thesis and through submissions to venues in empirical software 
engineering and process mining (e.g., EASE, MSR, ICPM).

Figure~\ref{fig:dsr-overview} illustrates the DSRM process applied to 
this research.

\begin{figure}[htbp]
\centering
\resizebox{\textwidth}{!}{%
\begin{tikzpicture}[
    node distance=0.3cm,
    phase/.style={rectangle, rounded corners=4pt, draw=black!70, fill=blue!8,
        text width=2.4cm, minimum height=3.0cm, align=center, font=\scriptsize},
    phaselabel/.style={font=\scriptsize\bfseries, text=black!80, 
        text width=2.4cm, align=center},
    artifact/.style={rectangle, rounded corners=2pt, draw=black!40, fill=gray!5,
        text width=2.2cm, minimum height=0.8cm, align=center, font=\tiny},
    arrow/.style={-{Stealth[length=4pt]}, thick, draw=black!50},
    eval/.style={rectangle, rounded corners=4pt, draw=black!70, fill=orange!12,
        text width=1.8cm, minimum height=1.0cm, align=center, font=\tiny\bfseries}
]

% Phase boxes
\node[phase] (ph1) {Problem\\Identification\\[6pt]\tiny Literature analysis\\3 problems\\Research gap};
\node[phaselabel, above=0.1cm of ph1] {Activity 1};

\node[phase, right=0.5cm of ph1] (ph2) {Objectives\\Definition\\[6pt]\tiny 3 RQs\\4 specific obj.\\5 contributions};
\node[phaselabel, above=0.1cm of ph2] {Activity 2};

\node[phase, right=0.5cm of ph2] (ph3) {Design \&\\Development\\[6pt]\tiny PATHCAST\\4-stage pipeline\\Python + PM4Py};
\node[phaselabel, above=0.1cm of ph3] {Activity 3};

\node[phase, right=0.5cm of ph3] (ph4) {Demonstration\\[12pt]\tiny 48 repositories\\227K commits\\Event logs};
\node[phaselabel, above=0.1cm of ph4] {Activity 4};

\node[phase, right=0.5cm of ph4] (ph5) {Evaluation\\[12pt]\tiny Two stages:\\E1, E2};
\node[phaselabel, above=0.1cm of ph5] {Activity 5};

\node[phase, right=0.5cm of ph5] (ph6) {Communication\\[12pt]\tiny Thesis +\\2 papers\\EASE, MSR};
\node[phaselabel, above=0.1cm of ph6] {Activity 6};

% Arrows
\draw[arrow] (ph1) -- (ph2);
\draw[arrow] (ph2) -- (ph3);
\draw[arrow] (ph3) -- (ph4);
\draw[arrow] (ph4) -- (ph5);
\draw[arrow] (ph5) -- (ph6);

% Evaluation stages
\node[eval, below=0.8cm of ph4.south east, anchor=north] (e1) {E1\\Pipeline\\Validation\\(RQ1)};
\node[eval, right=0.3cm of e1] (e2) {E2\\Forecasting\\Accuracy\\(RQ2, RQ3)};

\draw[arrow, dashed] (ph5.south) -- ++(0,-0.35) -| (e1.north);
\draw[arrow, dashed] (ph5.south) -- ++(0,-0.35) -| (e2.north);

% Artifacts below phases 1-3
\node[artifact, below=0.8cm of ph1] (a1) {Gap Analysis\\(Ch.~1, Ch.~2)};
\node[artifact, below=0.8cm of ph2] (a2) {RQs, Objectives\\Contributions};
\node[artifact, below=0.8cm of ph3] (a3) {PATHCAST\\Method (Ch.~4)};

\draw[arrow, dashed] (ph1.south) -- (a1.north);
\draw[arrow, dashed] (ph2.south) -- (a2.north);
\draw[arrow, dashed] (ph3.south) -- (a3.north);

\end{tikzpicture}%
}
\caption{Overview of the DSRM process applied to this research, with six 
activities, produced artifacts, and a two-stage evaluation.}
\label{fig:dsr-overview}
\end{figure}

The methodology is detailed in Chapter~\ref{ch:methodology}.


% =============================================================================
\section{Structure of the Document}
\label{sec:structure}

The remainder of this document is organized as follows.

Chapter~\ref{ch:theoretical-foundation} presents the theoretical
foundations, covering process mining, mining of software repositories,
Markov chains for software process modeling, Monte Carlo simulation for
process forecasting, and machine learning for software engineering.

Chapter~\ref{ch:slr} presents a systematic literature review following
Kitchenham and Charters' guidelines, mapping the state of the art at the
intersection of process mining, stochastic modeling, and machine learning applied to software process forecasting, and identifying the research gap addressed by the proposed method.

Chapter~\ref{ch:method} presents the proposed method (PATHCAST),
formalizing the four-stage pipeline with input/output contracts,
algorithms, formal properties, and theoretical guarantees.

Chapter~\ref{ch:methodology} details the research methodology, including
the Design Science Research framework, the experimental design for the
two evaluation stages, and threats to validity.

A dedicated implementation chapter (in preparation) will describe the implementation architecture, technology choices, and design decisions.

The empirical-evaluation chapter (in preparation) will present the empirical evaluation, including
the results of pipeline validation (E1) and forecasting accuracy assessment (E2), with statistical analysis and comparison against baselines.

The concluding chapter (in preparation) will discuss the results, identify limitations,
propose directions for future work, and present the conclusions.